\section{Local branches and global first-order deformations}
\label{sec:deformation}

\subsection{The analytic setting and the local lattices}
Throughout Sections~\ref{sec:deformation}--\ref{sec04:smoothing}, $X$
satisfies the projective Calabi--Yau hypotheses of the introduction.
All its singularities are isolated nodes or the exact anticanonical
cones over $E=\dP_5$, $\dP_6$, $\dP_7$, or $\PP^1\times\PP^1$.
The degree-five, degree-six and degree-seven surfaces are unique up
to isomorphism over $\CC$.  The degree-five surface is not toric;
the other three surfaces are toric.  Fix a nowhere-vanishing
section $\Omega$ of $\omega_X$.

At a node choose a small resolution; at a cone use the zero-section
contraction $\operatorname{Tot}(K_E)\to\Cone(E)$.  These proper analytic
maps are isomorphisms off the singular points and glue to a smooth
crepant analytic resolution $f_0:Y_0\to X$.  The gluing is Hausdorff
because two points over the same singular point lie in one proper
local model; points with distinct images can be separated on $X$.
Properness is local on the target.  Thus $Y_0$ is compact, and it is
Moishezon because it is bimeromorphic to the projective variety $X$.
The $\partial\bar\partial$-lemma holds on $Y_0$
\cite[Thm.~5.22 and Cor.~5.23]{dgms1975}.

Let $L_p$ be a small link of a singularity.  At a cone point the
circle-bundle Gysin sequence identifies
\[
 H_2(L_p;\QQ)=K_{E_p}^{\perp}\subset H_2(E_p;\QQ).
\]
For $E=\mathrm{Bl}_r\PP^2$, use the usual marking $h,e_1,\ldots,e_r$.
Its intersection form and canonical class are
\[
 \operatorname{diag}(1,-1,\ldots,-1),\qquad K_E=-3h+\sum_i e_i.
\]
At degree seven the selected line is
$\QQ(e_1-e_2)$.  For $E=\PP^1\times\PP^1$, the ruling classes
satisfy $f_1^2=f_2^2=0$, $f_1\cdot f_2=1$, and
$K_E=-2f_1-2f_2$; the selected line is the entire space
$K_E^\perp=\QQ(f_1-f_2)$.  At degree six set
\begin{equation}\label{eq:rootcoordinates}
 \ell=-h+e_1+e_2+e_3,\qquad A=-e_1+e_3,\qquad B=e_2-e_3.
\end{equation}
Then $K_E^\perp=\QQ\ell\oplus\langle A,B\rangle_\QQ$, the orthogonal
decomposition into the $A_1$ and $A_2$ root subspaces.  In the plane,
\begin{equation}\label{eq:a2coefficients}
 aA+bB=-a e_1+b e_2+(a-b)e_3.
\end{equation}
At a node the chosen resolution identifies its link-homology line
with the span of the exceptional curve.  Changing the small resolution
changes its generator by a sign and does not change any nonvanishing test.

At degree five the selected subspace is all of $K_E^\perp$.
For $E=\operatorname{Bl}_4\PP^2$ its five conic-pencil classes are
\begin{equation}\label{eq:degree5pencils}
 f_i=h-e_i\quad(1\le i\le4),\qquad
 f_5=2h-e_1-e_2-e_3-e_4,\qquad \sum_{i=1}^5f_i=-2K_E.
\end{equation}
For $\gamma=xh+\sum_i y_i e_i\in K_E^\perp$, define
\begin{equation}\label{eq:degree5coordinates}
 \lambda_i(\gamma)=\gamma\cdot f_i=x+y_i\quad(i\le4),
 \qquad \lambda_5(\gamma)=-x,\qquad \sum_i\lambda_i=0.
\end{equation}
These coordinates identify the integral perpendicular lattice with
the negative $A_4$ lattice: the inverse is
$x=-\lambda_5$, $y_i=\lambda_i+\lambda_5$, and
$\gamma^2=-\sum_i\lambda_i^2$.

\begin{definition}\label{def:branchprofile}
A branch profile $S$ selects the entire line at each node and each
$\PP^1\times\PP^1$ cone, the entire four-dimensional base at
each degree-five point, the unique reduced
degree-seven branch, and one of the two reduced branches at every
degree-six point.  Write $R_{p,S_p}$ for the corresponding subspace of
$H_2(L_p;\QQ)$ described above, and $K_S$ for \eqref{eq:relationkernel}.
We also write
\begin{equation}\label{eq:fullrelationkernel}
 K=\ker\left\{\bigoplus_pH_2(L_p;\QQ)\longrightarrow H_2(Y_0;\QQ)\right\}.
\end{equation}
\end{definition}

These subspaces are also
$R_{p,S_p}=\ker\{H_2(L_p;\QQ)\to H_2(\Mil_{p,S_p};\QQ)\}$, where
$\Mil_{p,S_p}$ is a smooth Milnor fibre on the chosen component
\cite[Prop.~2.1 and Lem.~4.2]{PaperIV}.  For the
$\PP^1\times\PP^1$ and degree-five cones the same assertion follows
from $H_2(\Mil_{p,S_p};\QQ)=0$, proved in
Lemmas~\ref{lem03:quadricmilnor} and~\ref{lem03:degree5milnor}.
When a profile is fixed, we
occasionally abbreviate $R_{p,S_p}$ to $\Rsm_p$.

The scheme-theoretic local bases are, in analytic coordinates,
\begin{equation}\label{eq:localbases}
 \begin{array}{c|c}
 \text{germ}&\text{local base ring}\\ \hline
 \text{node}&\CC\{t\}\\
 \Cone(\PP^1\times\PP^1)&\CC\{t\}\\
 \Cone(\dP_5)&\CC\{b_1,b_2,b_3,b_4\}\\
 \Cone(\dP_7)&\CC\{s,w\}/(w^2,sw)\\
 \Cone(\dP_6)&\CC\{s,u,v\}/(su,sv).
 \end{array}
\end{equation}
The degree-five entry is proved in
Proposition~\ref{prop03:degree5family}; the other cone entries are
Altmann's calculations.  At degree six the $A_1$ branch is $u=v=0$ and the $A_2$ branch is
$s=0$ \cite[\S\S5--8]{Altmann97}.  Fix analytic classifying maps from
a semiuniversal base $\Def(X)$ to these local bases.  Their
scheme-theoretic inverse image of the selected components is denoted
$\Def^S(X)$; the chosen maps form part of this notation.  Write
$V_S=T_0\Def^S(X)$.  The line and plane are smooth, but the plane's
smoothing locus excludes the three lines
$ab(a-b)=0$ in the marking \eqref{eq:a2coefficients}.  In particular,
the condition of being a nonzero tangent in the plane is weaker than
being a first-order smoothing direction.
At degree five the smoothing locus excludes the five hyperplanes
$\lambda_i=0$, as proved in Lemma~\ref{lem03:degree5marking}.

\subsection{The full local Hodge comparison}
Write $\mathbb T^i_X=\operatorname{Ext}^i(L_X,\cO_X)$, where $L_X$ is
the cotangent complex.  The following comparison permits the use of
K\"ahler differentials in the low-degree local-to-global sequence.

\begin{lemma}\label{lem:omegaL}
Let $X$ be a normal Gorenstein threefold whose non-lci locus is finite.  Then
\[
  \mathrm{Ext}^i(\Omega^1_X,\cO_X)\;\cong\;\mathrm{Ext}^i(L_X,\cO_X)
  \qquad(i\le2),
\]
both as sheaves and globally, compatibly with the local-to-global spectral
sequence.
\end{lemma}

\begin{proof}
The cohomology sheaves $\mathcal H^{-q}(L_X)$ for $q\ge1$ vanish where $X$ is a
local complete intersection, so they are supported on a finite set and have
finite length.  A finite-length module on a Cohen--Macaulay threefold has grade
$3$, so $\mathcal Ext^p(-,\cO_X)$ and $\mathrm{Ext}^p(-,\cO_X)$ vanish on it for
$p<3$.  In the spectral sequence
$E_2^{p,q}=\mathrm{Ext}^p(\mathcal H^{-q}(L_X),\cO_X)\Rightarrow
\mathrm{Ext}^{p+q}(L_X,\cO_X)$ a term with $q\ge1$ contributing in total degree
$\le2$ has $p\le1<3$, hence vanishes.  Only $q=0$ survives in degrees $\le2$, and
$\mathcal H^0(L_X)=\Omega^1_X$.
\end{proof}

The local-to-global sequence begins
\[
 0\longrightarrow H^1(X,\Theta_X)\longrightarrow\mathbb T^1_X
 \xrightarrow{\alpha}\bigoplus_pT^1_{X,p}
 \xrightarrow{\mathrm{ob}}H^2(X,\Theta_X).
\]
We use contraction with the restriction of the fixed global volume
form $\Omega$ in every local comparison below. Thus numerical local
parameters are Hodge-normalized; independently chosen equation
coordinates may differ by nonzero scalar factors.

\begin{lemma}[local Hodge comparison]\label{lem:localhodge}
Let $(Y,p)$ be a three-dimensional ordinary double point or the anticanonical
cone over $E=\dP_5$, $\dP_6$, $\dP_7$, or $\PP^1\times\PP^1$.  Let
$\pi_p\colon\widehat Y_p\to Y$ be the small resolution at a node and
$\widehat Y_p=\operatorname{Tot}(K_E)$ at a cone point, let $Z_p$ be its
exceptional locus, and put $U_p=\widehat Y_p\smallsetminus Z_p$.  Then the
connecting homomorphism gives an isomorphism
\begin{equation}\label{eq:localhodge}
 T^1_{Y,p}\;\xrightarrow{\ \sim\ }\;
 K_p^{\mathrm{cr}}:=
 \ker\!\left\{H^2_{Z_p}(\widehat Y_p;\Omega^2_{\widehat Y_p})
 \longrightarrow H^2(\widehat Y_p;\Omega^2_{\widehat Y_p})\right\}.
\end{equation}
There is also a natural isomorphism
\begin{equation}\label{eq:linklocalhodge}
 H_2(L_p;\CC)\cong H^3(L_p;\CC)
 \xrightarrow{\ \sim\ }K_p^{\mathrm{cr}},
\end{equation}
and the two isomorphisms are compatible with the local Gysin map to any global
crepant resolution containing $\widehat Y_p$.
\end{lemma}

\begin{proof}
The identifications
$T^1_{Y,p}\cong H^1(U_p;\Theta_{U_p})\cong H^1(U_p;\Omega^2_{U_p})$
are Schlessinger's comparison and contraction with the prescribed crepant volume
form.  The comparison applies to the non-lci cones: the hypothesis in
\cite[Lem.~1.17 and Lem.~4.1]{FL22a} is weakened in
\cite[Rem.~4.2]{FL22a} to an isolated singularity of depth at least $3$.

For a node,
$\widehat Y_p=\operatorname{Tot}_{\PP^1}(\cO(-1)\oplus\cO(-1))$.
The projection is affine, and the tangent sequence together with
$H^1(\PP^1,\cO(m))=0$ for $m\ge-1$ gives
$H^1(\widehat Y_p;\Theta_{\widehat Y_p})=0$.  The local cohomology sequence
therefore identifies the one-dimensional space $T^1_{Y,p}$ with
$K_p^{\mathrm{cr}}$.  The residue calculation for the node identifies this
class with the generator of $H^3(L_p)$; equivalently, this is the local
Gysin identification in Friedman's construction \cite{Friedman86}.

Now suppose $\widehat Y_p=\operatorname{Tot}(K_E)$, with projection
$\rho\colon\widehat Y_p\to E$ and zero section again denoted $E$.  Contraction
with the crepant volume form identifies
\[
 \Omega^2_{\widehat Y_p}(\log E)(-E)
 \cong \Theta_{\widehat Y_p}(-\log E),
\]
and the logarithmic tangent sequence is
\[
 0\longrightarrow\cO_{\widehat Y_p}\longrightarrow
 \Theta_{\widehat Y_p}(-\log E)\longrightarrow\rho^*\Theta_E
 \longrightarrow0.
\]
Since $\rho$ is affine and
$\rho_*\cO_{\widehat Y_p}=\bigoplus_{k\ge0}\cO_E(-kK_E)$, Kodaira vanishing
and Bott vanishing give
\begin{align*}
 H^1(\widehat Y_p;\cO_{\widehat Y_p})
   &=\bigoplus_{k\ge0}H^1(E;\cO_E(-kK_E))=0,\\
 H^1(\widehat Y_p;\rho^*\Theta_E)
   &=\bigoplus_{k\ge0}H^1
      \bigl(E;\Omega^1_E\otimes\cO_E(-(k+1)K_E)\bigr)=0.
\end{align*}
For the degree-five surface, the Bott vanishing used here is
Totaro's theorem \cite[Theorem~2.1]{totaro2020}; the twist
$-(k+1)K_E$ is ample for every $k\ge0$.  For the other surfaces
it is the usual toric Bott vanishing.
Thus
$H^1(\widehat Y_p;\Omega^2(\log E)(-E))=0$.  The exact sequence of
\cite[Thm.~2.1(v)]{FL22a}, whose theorem explicitly allows every
three-dimensional isolated canonical Gorenstein germ, now gives
\eqref{eq:localhodge}.  Moreover \cite[Thm.~2.1(vi)]{FL22a} gives a natural
inclusion
\[
 H^3(L_p;\CC)=\operatorname{Gr}^2_FH^3(L_p;\CC)
 \lhook\joinrel\longrightarrow K_p^{\mathrm{cr}}
\]
and a decomposition of $K_p^{\mathrm{cr}}$ with this as one summand.
If $E$ is toric and its boundary polygon has $n$ vertices, Altmann's calculation gives
$\dim T^1_{Y,p}=n-3$.  The circle-bundle Gysin sequence gives
$\dim H^3(L_p)=9-\deg E=n-3$.  Hence the other summand is zero and the displayed
inclusion is an isomorphism.  At degree five the same dimension
comparison uses $\dim T^1=4$ \cite[\S5.5]{Gross97rev} and
$\dim H^3(L_p)=9-5=4$.  Its construction in local cohomology is
compatible with the Gysin map, which proves the final assertion.
\end{proof}

\begin{lemma}[the local branches]\label{lem:branchhodge}
Under the isomorphism
$T^1_{Y,p}\cong H_2(L_p;\CC)$ of Lemma~\ref{lem:localhodge}, the tangent space
to each reduced local component is exactly the corresponding
subspace $\Rsm_p$: the whole line at a node or a
$\PP^1\times\PP^1$ cone, the entire $K_E^\perp$ at a degree-five
cone, the unique root line at a
$\dP_7$ cone, and respectively the $A_1$ line and the $A_2$ plane at a
$\dP_6$ cone.
\end{lemma}

\begin{proof}
At $E=\PP^1\times\PP^1$ both spaces are one-dimensional, and
$K_E^\perp=\CC(f_1-f_2)$ gives the assertion.  At degree five
the assertion follows from the full four-dimensional
comparison and the smooth irreducible base of
Proposition~\ref{prop03:degree5family}.  For the degree-six and
degree-seven points, consider a unit-edge polygon
$P=(v_0,\ldots,v_{n-1})$, put $e_i=v_{i+1}-v_i$.  Altmann identifies the
degree of the versal tangent space used here with
\begin{equation}\label{eq:altmanntangent}
 \mathcal M(P)=
 \left\{(t_i)\in\CC^n:\ \sum_i t_i e_i=0\right\}
 \big/\CC(1,\ldots,1),
\end{equation}
and a lattice Minkowski decomposition gives the class of its edge-dilation
vector $(t_i)$.  The Hodge isomorphism of Lemma~\ref{lem:localhodge} is natural
under lattice automorphisms of $P$, with one necessary character: contraction
with the crepant volume form twists equivariance by the determinant of the
induced automorphism of the three-dimensional cone lattice.
This equivariance is computed with the standard toric volume form.
The restriction of $\Omega$ is $g$ times that form, for an analytic
unit $g$.  Altmann's concentration of $T^1$ in the single Gorenstein
degree implies $\mathfrak m_pT^1_{X,p}=0$.  Multiplication by $g$
therefore acts on $T^1_{X,p}$ through the scalar $g(p)$.
The comparison determined by $\Omega$ differs only by this nonzero
scalar, so the invariant subspaces identified below are unchanged.

There is one unit-edge reflexive pentagon up to $\operatorname{GL}_2(\ZZ)$.  Take
\[
 P_7=\bigl((-1,-1),(0,-1),(1,0),(0,1),(-1,0)\bigr).
\]
Its nontrivial decomposition has dilation vector
$t=(0,1,0,1,0)$, up to replacing $t$ by $1-t$, and spans the reduced versal
base.  The reflection $s(x,y)=(y,x)$ fixes $[t]\in\mathcal M(P_7)$ and has
determinant $-1$.  On
$H_2(L_p)=K_E^\perp$ it acts by $-1$ on the unique root line and by $+1$ on
the complementary line.  The determinant twist therefore forces the Hodge
image of $[t]$ to be the root line.

Likewise there is one unit-edge reflexive hexagon.  Take
\[
 P_6=\bigl((1,0),(1,1),(0,1),(-1,0),(-1,-1),(0,-1)\bigr).
\]
The determinant-one rotation $r(x,y)=(x-y,x)$ cyclically permutes its vertices.
On $\mathcal M(P_6)$ its characteristic polynomial is
$(\lambda+1)(\lambda^2+\lambda+1)$.  The $(-1)$-eigenspace is spanned by the
alternating dilation $(0,1,0,1,0,1)$, the triangle-plus-triangle decomposition
which gives Altmann's one-dimensional smoothing component; the complementary
plane is spanned by the coarsenings of the three-segment decomposition and is
the two-dimensional component.  On
$K_E^\perp\cong A_1\oplus A_2$, the same rotation is $-1$ on $A_1$ and a
Coxeter element, with characteristic polynomial $\lambda^2+\lambda+1$, on
$A_2$.  Equivariance now identifies the line with $A_1$ and the plane with
$A_2$.  The boundary classes and edge-dilation vectors above give finite
intersection and representation calculations; Appendix~\ref{app:cert}
records their exact verification.
\end{proof}

\subsection{The image of the global tangent space}
The use of the full local spaces in Lemma~\ref{lem:localhodge} is
essential: the comparison also accounts for tangent directions that
do not belong to a reduced local component.

\begin{theorem}\label{thm:D}
Under the Hodge identifications determined by $\Omega$, localization
gives exact sequences
\begin{align}
 0&\longrightarrow H^1(X,\Theta_X)\longrightarrow\mathbb T^1_X
                  \longrightarrow K\otimes_\QQ\CC\longrightarrow0,
                  \label{eq:fulltangent}\\
 0&\longrightarrow H^1(X,\Theta_X)\longrightarrow V_S
                  \longrightarrow K_S\otimes_\QQ\CC\longrightarrow0.
                  \label{eq:selectedtangent}
\end{align}
\end{theorem}

\begin{proof}
Let $E=\operatorname{Exc}(f_0)$ and $U=Y_0\setminus E=X_{\mathrm{reg}}$.
Excision and the local cohomology sequence give
\begin{equation}\label{eq:globalsupport}
 H^1(U;\Omega^2_U)\longrightarrow
 \bigoplus_pH^2_{E_p}(\widehat X_p;\Omega^2_{\widehat X_p})
 \longrightarrow H^2(Y_0;\Omega^2_{Y_0}).
\end{equation}
Here $\widehat X_p$ denotes the restriction of $Y_0$ over a sufficiently
small neighborhood of $p$.  The algebraic cone calculations in
Lemma~\ref{lem:localhodge} compute these analytic local groups as well:
the vanishing statements are equivalent to the corresponding
vanishing of completed higher direct images, by formal functions.
The analytic local-support comparison itself is that of
\cite[Thm.~2.1]{FL22a}.

Under Lemma~\ref{lem:localhodge}, the map from
$\bigoplus_pK_p^{\mathrm{cr}}$ to $H^2(Y_0;\Omega^2_{Y_0})$, followed by
the de Rham map, is the topological link map followed by Poincar\'e
duality.  In particular the square
\[
\begin{array}{ccc}
 \displaystyle\bigoplus_pH_2(L_p;\CC)&\longrightarrow&H_2(Y_0;\CC)\\
 \downarrow\cong&&\downarrow\operatorname{PD}\\
 \displaystyle\bigoplus_pK_p^{\mathrm{cr}}&\longrightarrow&H^4(Y_0;\CC)
\end{array}
\]
commutes.  The $\partial\bar\partial$-lemma makes
$H^2(Y_0;\Omega^2_{Y_0})\to H^4(Y_0;\CC)$ injective.  Hence an element
of the sum of local tangent spaces has zero image in
$H^2(Y_0;\Omega^2_{Y_0})$ if and only if its link class lies in
$K\otimes_\QQ\CC$.

Contraction with $\Omega$ and Schlessinger's depth-three comparison
identify
\[
 H^1(U;\Omega^2_U)\cong H^1(U;\Theta_U)\cong\mathbb T^1_X;
\]
see \cite[Rem.~4.2]{FL22a}.  Exactness of \eqref{eq:globalsupport}
therefore gives every element of $K\otimes\CC$ as the localization
of a global tangent.  Conversely, the localization of any global
tangent has zero supported Gysin image, hence lies in that kernel.
The kernel of localization is $H^1(X,\Theta_X)$ by the local-to-global
sequence.  This proves \eqref{eq:fulltangent}.  Taking the inverse
image of the chosen local tangent subspaces and applying
Lemma~\ref{lem:branchhodge} proves \eqref{eq:selectedtangent}.
\end{proof}

At an $A_2$ point, \eqref{eq:selectedtangent} can supply a nonzero
local vector on a discriminant line.  It asserts the existence of a
global first-order deformation with those local components; it makes
no smoothing assertion for such a vector.  The next two sections
determine which directions integrate and which profiles contain
smooth fibres.
